% భాగం: అంతఃప్రజ్ఞావాద తర్కం
% అధ్యాయం: అర్థవిచారం
% విభాగం: ప్రతిపాదనలు

\documentclass[../../../include/open-logic-section]{subfiles}

\begin{document}

\olfileid{int}{sem}{pro}

\olsection{ప్రతిపాదనలు}

\tetoken{సంబంధ నమూనా}{!!a{relational model}}~$\mModel{M}$లో \tetoken{సూత్రం}{!!a{formula}}~$!A$ సత్యమయ్యే లోకాల సమితికి కొన్నిసార్లు పేరు పెట్టడం ఉపయోగపడుతుంది. ఆ సమితిని~$!A$ నిర్వచించే \emph{ప్రతిపాదన} అని పిలుస్తాం. దాన్ని $\Prop{M}{!A}$గా రాస్తాం.

\begin{defn}
  $V$ ప్రమేయాన్ని ఆగమన పద్ధతిలో కింది విధంగా విస్తరించి $\Prop{M}{!A} \colon \Frm
  \to \Pow{W}$ ప్రమేయాన్ని పొందుతాం.
  \begin{enumerate}
  \item $\Prop{M}{\lfalse} = \emptyset$
  \item $\Prop{M}{p} = V(p)$
  \item $\Prop{M}{\lnot !B} = {}$
    \[
      \Setabs{w \in W}{\text{$Rww'$ అయ్యే ప్రతి $w'$కు, $w'
          \notin \Prop{M}{!B}$}}.
    \]
  \item $\Prop{M}{!B \land !C} = \Prop{M}{!B} \cap \Prop{M}{!C}$.
  \item $\Prop{M}{!B \lor !C} = \Prop{M}{!B} \cup \Prop{M}{!C}$.
  \item $\Prop{M}{!B \lif !C} = {}$
    \[
      \Setabs{w \in W}{\text{$Rww'$ అయ్యే ప్రతి $w'$కు, $w'
          \in \Prop{M}{!B}$ అయితే $w' \in \Prop{M}{!C}$}}.
    \]
  \end{enumerate}
\end{defn}

\begin{prop}
  \ollabel{prop:propositions-true}
  $\mSat{M}{!A}[w]$ అవుతుంది అప్పుడే, అప్పుడు మాత్రమే $w \in \Prop{M}{!A}$.
\end{prop}

\begin{proof}
  సాధన.
\end{proof}

\begin{prob}
  \olref[int][sem][pro]{prop:propositions-true}ను నిరూపించండి.
\end{prob}

\end{document}
